% Solutions to the exercises of Chapter 15.
\begin{solutions}
  \item At rest the accelerometer measures the specific force $f = R_{wb}\T(a_w - g_w) = R_{wb}\T(0, 0, g)$ (\eqref{eq:imu-model} without bias and noise). With
        $R_{wb} = R_x(10^\circ)$,
        \[
          f = R_x(10^\circ)\T\begin{bmatrix}0\\0\\g\end{bmatrix} = \begin{bmatrix}0\\ g\sin10^\circ\\ g\cos10^\circ\end{bmatrix} = \begin{bmatrix}0\\ 1.70\\ 9.66\end{bmatrix}~\mathrm{m/s^2}.
        \]
        The $y$-component reveals the tilt, with a positive sign: a positive rotation about $x$ lifts the body $y$-axis ($R_{wb}e_y = (0, \cos10^\circ, \sin10^\circ)$), so the
        upward reaction to gravity has a positive component along it. The magnitude is still $g$; only its direction in the body frame has changed.

  \item At frequencies far below resonance the deflection is $x = (a - g)/\omega_n^2$ (\cref{eq:accel-dynamics}), so the static sensitivity is
        \[
          \frac{1}{\omega_n^2} = \frac{1}{(2\pi\cdot5000~\mathrm{Hz})^2} = 1.0\times10^{-9}~\mathrm{m\ per\ m/s^2},
        \]
        one nanometre per $\mathrm{m/s^2}$, about five silicon atoms; the full $1g$ deflects the proof mass by about 10~nm. A noise floor of
        $2\times10^{-3}~\mathrm{m/s^2/\sqrt{Hz}}$ over a 100~Hz bandwidth ($0.02~\mathrm{m/s^2}$) corresponds to a displacement of $2\times10^{-11}$~m, a tenth of an atom.
        This is possible only because the capacitive readout measures the \emph{average} position of a large electrode, over many atoms, at a
        frequency where thermal and electronic noise can be filtered. It also explains the trade-off of \cref{ch:imu}: a stiffer (higher-frequency)
        sensor has more bandwidth but smaller deflections and therefore more relative noise.

  \item In the frame of the sensor, which rotates with angular velocity $\Omega$, \eqref{eq:coriolis} says that a mass moving with relative velocity $\dot r$
        experiences, besides its true acceleration, the apparent accelerations $-2\,\Omega\times\dot r$ and $-\Omega\times(\Omega\times r)$. For the drive motion
        $\dot r = A\omega_d\cos(\omega_dt)\,e_x$ and a rate $\Omega_z$ about $z$,
        \[
          -2\,\Omega\times\dot r = -2\Omega_zA\omega_d\cos(\omega_dt)\,(e_z\times e_x) = -2\Omega_zA\omega_d\cos(\omega_dt)\,e_y ,
        \]
        which is \eqref{eq:gyro-coriolis}: an oscillation along $y$ at the drive frequency, with amplitude proportional to $\Omega_z$. The centripetal term is quadratic in $\Omega$ and does not oscillate at $\omega_d$, so demodulation at the drive frequency rejects it.
        In a tuning fork the two masses are driven in anti-phase: their velocities, and therefore their Coriolis accelerations, are opposite, while a
        linear acceleration pushes both the same way. The differential readout (difference of the two sense displacements) doubles the Coriolis
        signal and cancels the common linear acceleration; imperfect matching of the two masses leaves the small residual $g$-sensitivity of
        \cref{tab:imu-errors}.

  \item With the formulas of the six-position test,
        \[
          b = \frac{9.83 + (-9.79)}{2} = 0.02~\mathrm{m/s^2} \ (\approx 2~\mathrm{mg}), \qquad 1 + s = \frac{9.83 - (-9.79)}{2\cdot9.81} = 1.000 ,
        \]
        so this axis has a bias of $0.02~\mathrm{m/s^2}$ and no measurable scale error. (The scale factor is relative to the gravity value used; with the local
        $g$ instead of $9.81$, $s$ would change by the relative difference, a few parts in $10^4$ in practice.)

  \item Using $1\,^\circ/\sqrt{\mathrm h} = \frac{\pi/180~\mathrm{rad}}{\sqrt{3600~\mathrm s}} = 2.909\times10^{-4}~\mathrm{rad/s/\sqrt{Hz}}$, an ARW of $0.2\,^\circ/\sqrt{\mathrm h}$ is
        $\sigma_g = 5.82\times10^{-5}~\mathrm{rad/s/\sqrt{Hz}}$. At 400~Hz each sample has standard deviation $\sigma_g/\sqrt\Dt = 5.82\times10^{-5}\sqrt{400} =
        1.16\times10^{-3}~\mathrm{rad/s}$. For comparison, EuRoC's gyroscope ($1.7\times10^{-4}$) has about three times this ARW.

  \item The reading is $(1 + s)\omega$, so the error is $s\omega = 10^{-3}~\mathrm{rad/s}$ at 1~rad/s: an equivalent bias of $10^{-3}$~rad/s ($0.06\,^\circ$/s), larger than the
        EuRoC gyro's noise over most horizons. A constant-bias model cannot absorb it during maneuvers because the error is proportional to the rate:
        it changes sign when the vehicle turns the other way, vanishes when it stops, and doubles when it turns twice as fast. A random-walk bias
        changes slowly by assumption, so the filter either lags behind these changes or, if the bias process noise is inflated to follow them,
        loses its ability to estimate the true bias. Each full turn accumulates an attitude error of $2\pi s = 6.3$~mrad, in the plane of the turn. The
        remedies are calibration (on a rate table) or estimating $s$ as an additional state when the vehicle rotates enough to make it observable.
\end{solutions}
